\subsection{Range} 
We find the range of a piecewise function in the same way we find the domain, but it can be more challenging, because the ranges of the pieces can \makeblank[1in]. They may also be harder to find from the expressions.
\begin{Example}
Give the domain of $f(x)=\pw{3x&-2&\leq&x&<&1\\4-x&1&<&x&<&5\\6&5&\leq&x&\leq&6\\-2x+4&x&>&8}$
\begin{lpic}{}{12}
\end{lpic}
$\rng f=\makeblank$
\end{Example}
If we are looking at the graph of a piecewise function, we go back to our most basic technique: \makeblank[1in] the graph into the \makeblanksmall axis.
\begin{Example}
Give the domain of the piecewise function graphed below.
\begin{icpic}[x=0.5\linespace,y=0.5\linespace]
\setgraphsquare{5}
\setspace{1}
\GraphSmall
\draw[graphend] (-3,4) -- (-2,3);
\draw[graphend] (-2,-3) -- (1,2);
\draw[graphend] (2,-4) -- (3,-2);
\filldraw[dotempty] (-3,4) circle (0.2);
\filldraw[dotempty] (-2,3) circle (0.2);
\filldraw[dotfull] (-2,-3) circle (0.2);
\filldraw[dotfull] (3,-2) circle (0.2);
\filldraw[dotempty] (1,2) circle (0.2);
\filldraw[dotempty] (2,-4) circle (0.2);
\end{icpic}
$\rng f=\makeblank$
\end{Example}